\documentclass[10pt,a4paper]{amsart}
\usepackage{amsmath,amssymb,amsthm,graphicx}
% For commutative diagrams we use Xy-pic
\usepackage[all]{xy}

% We use 2cell for 2-commutative diagrams.
\xyoption{2cell}
% Theorem environments.
%
\theoremstyle{plain}
\newtheorem{theorem}[subsection]{Theorem}
\newtheorem{proposition}[subsection]{Proposition}
\newtheorem{lemma}[subsection]{Lemma}

\theoremstyle{definition}
\newtheorem{definition}[subsection]{Definition}
\newtheorem{example}[subsection]{Example}
\newtheorem{exercise}[subsection]{Exercise}
\newtheorem{situation}[subsection]{Situation}

\theoremstyle{remark}
\newtheorem{remark}[subsection]{Remark}
\newtheorem{remarks}[subsection]{Remarks}

\numberwithin{equation}{subsection}

% Macros
%
\def\lim{\mathop{\mathrm{lim}}\nolimits}
\def\colim{\mathop{\mathrm{colim}}\nolimits}
\def\Spec{\mathop{\mathrm{Spec}}}
\def\Hom{\mathop{\mathrm{Hom}}\nolimits}
\def\Ext{\mathop{\mathrm{Ext}}\nolimits}
\def\SheafHom{\mathop{\mathcal{H}\!\mathit{om}}\nolimits}
\def\SheafExt{\mathop{\mathcal{E}\!\mathit{xt}}\nolimits}
\def\Sch{\mathit{Sch}}
\def\Mor{\mathop{\mathrm{Mor}}\nolimits}
\def\Ob{\mathop{\mathrm{Ob}}\nolimits}
\def\Sh{\mathop{\mathit{Sh}}\nolimits}
\def\NL{\mathop{N\!L}\nolimits}
\def\CH{\mathop{\mathrm{CH}}\nolimits}
\def\proetale{{pro\text{-}\acute{e}tale}}
\def\etale{{\acute{e}tale}}
\def\QCoh{\mathit{QCoh}}
\def\Ker{\mathop{\mathrm{Ker}}}
\def\Im{\mathop{\mathrm{Im}}}
\def\Coker{\mathop{\mathrm{Coker}}}
\def\Coim{\mathop{\mathrm{Coim}}}

% Boxtimes
%
\DeclareMathSymbol{\boxtimes}{\mathbin}{AMSa}{"02}

%
% Macros for moduli stacks/spaces
%
\def\QCohstack{\mathcal{QC}\!\mathit{oh}}
\def\Cohstack{\mathcal{C}\!\mathit{oh}}
\def\Spacesstack{\mathcal{S}\!\mathit{paces}}
\def\Quotfunctor{\mathrm{Quot}}
\def\Hilbfunctor{\mathrm{Hilb}}
\def\Curvesstack{\mathcal{C}\!\mathit{urves}}
\def\Polarizedstack{\mathcal{P}\!\mathit{olarized}}
\def\Complexesstack{\mathcal{C}\!\mathit{omplexes}}
% \Pic is the operator that assigns to X its picard group, usage \Pic(X)
% \Picardstack_{X/B} denotes the Picard stack of X over B
% \Picardfunctor_{X/B} denotes the Picard functor of X over B
\def\Pic{\mathop{\mathrm{Pic}}\nolimits}
\def\Picardstack{\mathcal{P}\!\mathit{ic}}
\def\Picardfunctor{\mathrm{Pic}}
\def\Deformationcategory{\mathcal{D}\!\mathit{ef}}
\usepackage{cleveref}
\crefname{originallemma}{original lemma}{original lemmas}
\Crefname{originallemma}{Original lemma}{Original lemmas}
\makeatletter
\newlabel{lemma-additivity-periodic-length}{{0EA7}{source}{}{sourcepointer-1}{}}
\newlabel{lemma-additivity-periodic-length@cref}{{[originallemma][0][]0EA7}{[1][0][]source}{}{}{}}
\makeatother
\title{The multiplicity of a two-periodic complex of finite length modules is the alternating sum of the lengths}
\author{The Stacks Project Authors}
\date{stacks-project:chow.tex:local-snapshot-20261011; GNU FDL 1.2 or later}
\begin{document}
\maketitle
\noindent\textit{Editorial note.} All delivered mathematical text is the author's own, copied from the pinned Stacks Project chapter file chow.tex; no statement or proof step was rewritten, reordered or removed. The author leaves the two reduced cases to the reader; that sentence is part of the author's text and is reproduced as written.\par\medskip
\noindent\textit{Editorial note.} The preamble is the chapter's own native preamble excerpt used for this source: the Xy-pic setup at lines 14-18 and the theorem declarations with the mathematical macros at lines 154-217 of sources/stacks-chow/preamble.tex. The quoted class line and the xr-hyper externaldocument block are not included because the assembler supplies a plain amsart wrapper; no macro was invented or redefined.\par\medskip
\noindent\textit{Editorial note.} The same-chapter referenced lemma is located at its \textbackslash{}label line inside the pinned chapter file.\par\medskip
\noindent\textit{Allowed standard prerequisites (editorial).}\begin{itemize}
\item Two-periodic complexes and their multiplicity e\_R. Same chapter, Definition environment labelled definition-periodic-complex (source line 110) together with the displayed definition of e\_R at source lines 200-205 (label equation-multiplicity-coker-ker, line 203). Supplies e\_R(M, N, phi, psi) and the (2,1)-periodic specialisation e\_R(M, phi, psi); the reduction to the two extreme cases is the counted content.
\item Additivity of the multiplicity in short exact sequences. Same chapter, Lemma labelled lemma-additivity-periodic-length (source line 209, Stacks tag 0EA7); invoked at proof line 277. Supplies the additivity input used to cut the problem down to the two special cases; the explicit short exact sequence and the additivity of length(M) - length(N) are inside the counted span.
\end{itemize}\medskip
\ifdefined\hypertarget\hypertarget{sourcepointer-1}{}\fi
\noindent\textit{Source reference.} Original lemma 0EA7: sources/stacks-chow/pinned-original-source.tex, source lines 209--209. Same-chapter pointer invoked at proof line 277 to split the claim into the two cases M = 0 and phi surjective; the additivity remark, the displayed short exact sequence and the reduction are retained verbatim in the counted span.\par\medskip
\begin{lemma}
\label{lemma-finite-periodic-length}
Let $R$ be a ring. If $(M, N, \varphi, \psi)$ is a $2$-periodic complex
such that $M$, $N$ have finite length, then
$e_R(M, N, \varphi, \psi) = \text{length}_R(M) - \text{length}_R(N)$.
In particular, if $(M, \varphi, \psi)$ is a $(2, 1)$-periodic complex
such that $M$ has finite length, then
$e_R(M, \varphi, \psi) = 0$.
\end{lemma}

\begin{proof}
Observe that on the category of $2$-periodic complexes with $M$, $N$
of finite length the quantity ``$\text{length}_R(M) - \text{length}_R(N)$''
is additive in short exact sequences (precise statement left to the reader).
Consider the short exact sequence
$$
0 \to (M, \Im(\varphi), \varphi, 0) \to
(M, N, \varphi, \psi) \to (0, N/\Im(\varphi), 0, 0) \to 0
$$
The initial remark combined with the additivity of
Lemma \ref{lemma-additivity-periodic-length}
reduces us to the cases (a) $M = 0$ and (b) $\varphi$ is surjective.
We leave those cases to the reader.
\end{proof}

\end{document}
